% Unit 080 terjemahan Bahasa Indonesia dari chapter6.tex baris 1540--1631.
% Sumber: Methods of Algebra, Volume 2, commit
% 9a5803ff77dd3257484cb177f851a73770a59dd3 (CC BY 4.0).
% stable-unit-id: o014.aljabr2.chapter6.lhs-spectral-sequence-a
% Cakupan: Bagian 6.9, pembentukan dan pernyataan teorema LHS.

% segment-id: o014.li2.chapter6.lhs-spectral-sequence-a.h001
\section{Barisan spektral Lyndon--Hochschild--Serre}\label{sec:LHS-SS}
\hypertarget{o014.aljabr2.chapter6.lhs-spectral-sequence-a}{ }

% segment-id: o014.li2.chapter6.lhs-spectral-sequence-a.p001
Sepanjang bagian ini, gunakan gelanggang komutatif $\Bbbk$, grup $G$, dan
subgrup normal $H\lhd G$ yang tetap. Ingat kembali funktor restriksi
$\Res^G_H:G\dcate{Mod}\to H\dcate{Mod}$ dan funktor inflasi
$\mathrm{Infl}^G_{G/H}:G/H\dcate{Mod}\to G\dcate{Mod}$ dari
Definisi~\ref{def:Res-Infl}. Untuk setiap modul-$G$ $M$, perkenalkan notasi
\[ M^{H\lhd G}:=(\Res^G_HM)^H,\qquad
	M_{H\lhd G}:=(\Res^G_HM)_H. \]
Perhatikan bahwa $M^{H\lhd G}$ mempunyai aksi kiri $G/H$ yang jelas,
$(gH)\cdot x:=gx$. Dengan cara serupa, $M_{H\lhd G}$ mempunyai aksi kiri
$G/H$ yang jelas. Ini memberikan dua funktor
\[ (\cdot)^{H\lhd G},\ (\cdot)_{H\lhd G}:
	G\dcate{Mod}\to G/H\dcate{Mod}. \]

% segment-id: o014.li2.chapter6.lhs-spectral-sequence-a.l001
\begin{lemma}
	Funktor $(\cdot)^{H\lhd G}$ adalah adjoin kanan dari
	$\mathrm{Infl}^G_{G/H}$, sedangkan $(\cdot)_{H\lhd G}$ adalah adjoin
	kirinya.
\end{lemma}
\begin{proof}
	Ini merupakan perluasan sederhana dari
	Proposisi~\ref{prop:inv-coinv-infl}; verifikasinya diserahkan kepada
	pembaca.
\end{proof}

% segment-id: o014.li2.chapter6.lhs-spectral-sequence-a.p002
Jelas terdapat diagram komutatif funktor
\[\begin{tikzcd}[column sep=tiny]
	G\dcate{Mod} \arrow[rr,"{(\cdot)^{H\lhd G}}"]
	\arrow[rd,"{(\cdot)^G}"'] & &
	G/H\dcate{Mod} \arrow[ld,"{(\cdot)^{G/H}}"]\\
	& \Bbbk\dcate{Mod} &
\end{tikzcd}\quad\begin{tikzcd}[column sep=tiny]
	G\dcate{Mod} \arrow[rr,"{(\cdot)_{H\lhd G}}"]
	\arrow[rd,"{(\cdot)_G}"'] & &
	G/H\dcate{Mod}. \arrow[ld,"{(\cdot)_{G/H}}"]\\
	& \Bbbk\dcate{Mod} &
\end{tikzcd}\]

% segment-id: o014.li2.chapter6.lhs-spectral-sequence-a.p003
Funktor $(\cdot)^{H\lhd G}$ eksak kiri dan mempertahankan objek injektif,
sebab ia mempunyai adjoin kiri eksak $\mathrm{Infl}^G_{G/H}$. Dengan cara
yang sama, $(\cdot)_{H\lhd G}$ merupakan funktor eksak kanan yang
mempertahankan objek projektif. Maka Teorema~\ref{prop:derived-composite}
mengenai penurunan funktor komposit memberikan isomorfisme kanonik dalam
kategori turunan $\cate{D}(\Bbbk\dcate{Mod})$ (dengan tanda kurung yang sesuai
dihilangkan):
\begin{equation}\label{eqn:LHS-derived}
	\begin{gathered}
		\mathrm{R}(\cdot)^G\rightiso
		\mathrm{R}(\cdot)^{G/H}\circ\mathrm{R}(\cdot)^{H\lhd G},\\
		\mathrm{L}(\cdot)_G\leftiso
		\mathrm{L}(\cdot)_{G/H}\circ\mathrm{L}(\cdot)_{H\lhd G}.
	\end{gathered}
\end{equation}

% segment-id: o014.li2.chapter6.lhs-spectral-sequence-a.p004
Funktor turunan $\mathrm{R}(\cdot)^{H\lhd G}$ dan
$\mathrm{L}(\cdot)_{H\lhd G}$ dalam rumus tersebut bukanlah objek baru.
Untuk masing-masing diagram komutatif
\[\begin{tikzcd}
	G\dcate{Mod} \arrow[r,"{(\cdot)^{H\lhd G}}"]
	\arrow[d,"{\Res^G_H}"'] & G/H\dcate{Mod}
	\arrow[d,"{\text{pelupa}=\Res^{G/H}_{\{1\}}}"]\\
	H\dcate{Mod} \arrow[r,"{(\cdot)^H}"'] & \Bbbk\dcate{Mod}
\end{tikzcd}\quad\begin{tikzcd}
	G\dcate{Mod} \arrow[r,"{(\cdot)_{H\lhd G}}"]
	\arrow[d,"{\Res^G_H}"'] & G/H\dcate{Mod}
	\arrow[d,"\text{pelupa}"]\\
	H\dcate{Mod} \arrow[r,"{(\cdot)_H}"'] & \Bbbk\dcate{Mod},
\end{tikzcd}\]
turunkan kedua komposisi menurut Teorema~\ref{prop:derived-composite}. Karena
$\Res^G_H$ eksak dan mempertahankan objek injektif maupun projektif
(Proposisi~\ref{prop:Ind-preservation}), sedangkan funktor pelupa eksak,
teorema tersebut memberikan
\begin{gather*}
	\text{pelupa}\circ\mathrm{R}(\cdot)^{H\lhd G}
	\leftiso\mathrm{R}(\text{funktor komposit})
	\rightiso\mathrm{R}(\cdot)^H\circ\Res^G_H,\\
	\text{pelupa}\circ\mathrm{L}(\cdot)_{H\lhd G}
	\rightiso\mathrm{L}(\text{funktor komposit})
	\leftiso\mathrm{L}(\cdot)_H\circ\Res^G_H.
\end{gather*}
Artinya, setelah menerapkan funktor pelupa, funktor turunan kanan
$\mathrm{R}(\cdot)^{H\lhd G}$ sama dengan $\mathrm{R}(\cdot)^H$, dan funktor
turunan kiri $\mathrm{L}(\cdot)_{H\lhd G}$ sama dengan
$\mathrm{L}(\cdot)_H$. Karena itu, kohomologi dan homologinya wajar ditulis
sebagai
\begin{align*}
	\mathrm{R}^n(\cdot)^{H\lhd G}
	&=\Hm^n(H,\cdot)\quad\text{dengan aksi $G/H$},\\
	\mathrm{L}_n(\cdot)_{H\lhd G}
	&=\Hm_n(H,\cdot)\quad\text{dengan aksi $G/H$}.
\end{align*}
Perhatikan bahwa kita mulai tidak lagi mencantumkan funktor restriksi $\Res$
dalam notasi.

% segment-id: o014.li2.chapter6.lhs-spectral-sequence-a.p005
Aksi-aksi ini mudah dideskripsikan. Untuk kohomologi, misalkan $M$ modul-$G$
dan $n\in\Z_{\geq0}$. Setidaknya ada tiga sudut pandang bagi aksi $G/H$:
\begin{enumerate}
	\item Ambil resolusi injektif $0\to M\to I^0\to\cdots$. Resolusi ini juga
	merupakan resolusi injektif modul-$H$. Grup $G/H$ beraksi secara alami pada
	kompleks $\cdots\to I^{n,H}\to I^{n+1,H}\to\cdots$, dan karenanya pada
	$\Hm^n(H,M)$. Deskripsi ini langsung berasal dari konstruksi abstrak di
	atas.
	\item Grup $G/H$ beraksi pada funktor
	$(\cdot)^H:G\dcate{Mod}\to\Bbbk\dcate{Mod}$. Sifat universal funktor-$\delta$
	kohomologis universal kemudian membuatnya beraksi pada $\Hm^n(H,M)$,
	kompatibel dengan barisan eksak panjang.
	\item Dalam Contoh~\ref{eg:change-of-group-conj}, setiap $t\in G/H$ beraksi
	pada $\Hm^n(H,M)$ melalui $\Hm^n(c_t)$.
\end{enumerate}
Pembaca diminta memeriksa bahwa ketiga deskripsi ini sama. Intinya ialah
mereduksi perbandingan tersebut ke kasus $n=0$.

% segment-id: o014.li2.chapter6.lhs-spectral-sequence-a.p006
Hal penting berikutnya ialah mengekstrak informasi yang lebih konkret dari
\eqref{eqn:LHS-derived}. Ini melibatkan barisan spektral bergradasi ganda
kohomologis dan homologis dari Definisi~\ref{def:graded-spectral-sequence},
serta persiapan berikut.

% segment-id: o014.li2.chapter6.lhs-spectral-sequence-a.l002
\begin{lemma}\label{prop:LHS-action-aux}
	Untuk setiap modul-$G$ $M$ dan setiap $n$, homomorfisme kanonik
	$\Hm^n(G,M)\to\Hm^n(H,M)$ (lihat \eqref{eqn:change-of-group-coh}) berfaktor
	melalui $\Hm^n(H,M)^{G/H}$. Homomorfisme kanonik
	$\Hm_n(H,M)\to\Hm_n(G,M)$ berfaktor melalui
	$\Hm_n(H,M)_{G/H}$.
\end{lemma}
\begin{proof}
	Kita bahas kohomologi. Ambil resolusi injektif $M\to I$, dengan
	$I=[I^0\to I^1\to\cdots]$, dan tuliskan
	$I^H:=[I^{0,H}\to\cdots]$. Karena $I$ juga merupakan resolusi injektif dari
	$M$ sebagai modul-$H$, penyertaan $I^G\hookrightarrow I^H$ menginduksi
	$\Hm^n(G,M)\to\Hm^n(H,M)$.

	Seperti yang diharapkan, Catatan~\ref{rem:change-of-group-coh} memastikan
	bahwa pemetaan ini tepat merupakan homomorfisme kanonik yang dimaksud.
	Sekarang bandingkan dengan pembahasan aksi $G/H$ di atas.
\end{proof}

% segment-id: o014.li2.chapter6.lhs-spectral-sequence-a.t001
\begin{theorem}[R.\ Lyndon, G.\ Hochschild--J.-P.\ Serre]
\label{prop:LHS-SS}
	\index{barisan spektral!Lyndon--Hochschild--Serre}
	Misalkan $M$ modul-$G$ dan $H\lhd G$.
	\begin{enumerate}[(i)]
		\item Terdapat barisan spektral kohomologis dan homologis di kuadran
		pertama
		\begin{align*}
			E_2^{p,q}&\simeq\Hm^p\left(G/H,\Hm^q(H,M)\right)
			\Rightarrow\Hm^{p+q}(G,M),\\
			E^2_{p,q}&\simeq\Hm_p\left(G/H,\Hm_q(H,M)\right)
			\Rightarrow\Hm_{p+q}(G,M).
		\end{align*}
		\item Misalkan $n\in\Z_{\geq1}$ dan andaikan
		$\Hm^i(H,M)=0$ untuk semua $1\leq i<n$ dalam kasus kohomologis, atau
		$\Hm_i(H,M)=0$ untuk semua $1\leq i<n$ dalam kasus homologis. Maka
		masing-masing terdapat barisan eksak
		\begin{equation*}\begin{split}
			0\to\Hm^n(G/H,M^H)&\to\Hm^n(G,M)\to\Hm^n(H,M)^{G/H}\\
			&\to\Hm^{n+1}(G/H,M^H)\to\Hm^{n+1}(G,M),
		\end{split}\end{equation*}
		dan
		\begin{equation*}\begin{split}
			\Hm_{n+1}(G,M)&\to\Hm_{n+1}(G/H,M_H)\\
			&\to\Hm_n(H,M)_{G/H}\to\Hm_n(G,M)
			\to\Hm_n(G/H,M_H)\to0.
		\end{split}\end{equation*}
		Perhatikan bahwa untuk $n=1$, prasyarat ini berlaku secara hampa.
		\item Dalam barisan eksak kohomologis di atas, semua morfisme selain
		$\Hm^n(H,M)^{G/H}\to\Hm^{n+1}(G/H,M^H)$ berasal dari homomorfisme
		ekuivarian $M^H\hookrightarrow M$ (terhadap
		$G/H\twoheadleftarrow G$; lihat Definisi~\ref{def:change-of-group-equiv}),
		homomorfisme $M\xrightarrow{\identity}M$ (terhadap
		$G\hookleftarrow H$), serta faktorisasi dalam
		Lema~\ref{prop:LHS-action-aux}. Kasus barisan eksak homologis serupa.
	\end{enumerate}
\end{theorem}
